\documentclass[11pt]{article}
\usepackage[margin=1in]{geometry}
\usepackage{amsmath,amssymb,amsthm,mathtools}
\usepackage{microtype}
\usepackage[colorlinks=true,linkcolor=blue,citecolor=blue,urlcolor=blue]{hyperref}
\newtheorem{theorem}{Theorem}[section]
\newtheorem{corollary}[theorem]{Corollary}
\newtheorem{lemma}[theorem]{Lemma}
\newtheorem{proposition}[theorem]{Proposition}
\theoremstyle{remark}
\newtheorem{remark}[theorem]{Remark}
\newcommand{\R}{\mathbb R}
\newcommand{\Q}{\mathbb Q}
\newcommand{\E}{\mathbb E}
\newcommand{\Prob}{\mathbb P}
\newcommand{\one}{\mathbf 1}
\newcommand{\norm}[1]{\left\lVert#1\right\rVert}
\newcommand{\ip}[2]{\langle#1,#2\rangle}
\newcommand{\diag}{\operatorname{diag}}
\newcommand{\im}{\operatorname{im}}
\newcommand{\supp}{\operatorname{supp}}
\title{Vanishing of first path homology above the\texorpdfstring{\\}{ }random-digraph exponent $-2/3$}
\author{Kaiwen Jin\\\small\href{mailto:yui.ui.siki@gmail.com}{yui.ui.siki@gmail.com}}
\date{October 14, 2026}
\begin{document}
\maketitle
\begin{abstract}
Let $G(n,p)$ be the loop-free random digraph in which every ordered edge is
present independently with probability $p$, with reciprocal edges permitted.
We prove that its first non-regular path homology over $\Q$ vanishes with
probability tending to one if $n^2p^3/\log^4(2n)\to\infty$.
Consequently, vanishing holds for $p=n^\alpha$ for every fixed
$-2/3<\alpha\leq0$, as conjectured by Chaplin.
The proof uses only boundaries of directed diamonds and transitive triangles.
A uniform estimate for the diamond quadratic form shows that its annihilator
consists of vertex gradients and the constant edge cochain. A transitive
triangle removes the latter direction. The estimates retain the independent
ordered-edge model and the non-regular boundary convention throughout.
\end{abstract}

\section{Introduction}

Path homology records cycles in a directed graph using relations supplied by
allowed directed paths. For random directed graphs it can therefore distinguish
directed-network structure from an independent-edge null model.
Chaplin~\cite{Chaplin} studied the first path Betti number in this model and
proved its vanishing when $p\gg(\log n/n)^{1/3}$; see
Proposition~4.17 of that paper. Conjecture~B.1 predicts vanishing throughout the
larger power-law range $p=n^\alpha$, $\alpha>-2/3$.
The present paper proves this prediction for non-regular path homology over
$\Q$ by controlling all edge cochains at once.

We first specify the chain convention. For a digraph $G$ on $[n]=\{1,\ldots,n\}$,
let $\Lambda_k(\Q)$ have as basis all vertex sequences
$[v_0,\ldots,v_k]$, allowing repetitions, with boundary
\[
 \partial[v_0,\ldots,v_k]
 =\sum_{j=0}^k(-1)^j[v_0,\ldots,\widehat v_j,\ldots,v_k].
\]
Let $A_k$ be the subspace spanned by paths whose successive ordered edges
belong to $G$. Put $\Omega_0=A_0$ and
$\Omega_k=A_k\cap\partial^{-1}(A_{k-1})$ for $k\geq1$.
No boundary term with repeated consecutive vertices is discarded.
The first non-regular path Betti number over $\Q$ is
\[
 \beta_1^{\rm nr}(G)
 =\dim_\Q\frac{\ker(\partial:\Omega_1\to\Omega_0)}{\partial\Omega_2}.
\]
Since $G$ has no loops, $\Omega_1=A_1$ is the vector space on its ordered
edges. The identity $\partial^2=0$ gives the indicated quotient.

Write $G(n,p)$ for the digraph with adjacency variables $A_{ab}$,
$a\ne b$, independent Bernoulli$(p)$, and $A_{aa}=0$. In particular
$A_{ab}$ and $A_{ba}$ are independent. All asymptotics below concern
deterministic sequences of probabilities as $n\to\infty$.

\begin{theorem}\label{thm:main}
Suppose $0<p_n\leq1$ and
\begin{equation}\label{eq:condition}
 \frac{n^2p_n^3}{\log^4(2n)}\longrightarrow\infty.
\end{equation}
Then
$\Prob\{\beta_1^{\rm nr}(G(n,p_n))=0\}\longrightarrow1$.
\end{theorem}

\begin{corollary}\label{cor:power}
For every fixed $-2/3<\alpha\leq0$, a digraph with independent ordered
edges of probability $p_n=n^\alpha$ satisfies
$\Prob\{\beta_1^{\rm nr}(G(n,p_n))=0\}\to1$.
\end{corollary}
\begin{proof}
Here $n^2p_n^3=n^{2+3\alpha}$ has a strictly positive exponent and hence
dominates $\log^4(2n)$. Apply Theorem~\ref{thm:main}.
\end{proof}

The sufficient condition~\eqref{eq:condition} identifies the conjectured
power-law exponent without determining the optimal logarithmic factor.
Neither the critical exponent $\alpha=-2/3$ nor convergence uniform in
$\alpha$ is asserted. The conclusion concerns rational homology; it does
not assert the absence of integral torsion.

The mechanism is visible from two elementary boundaries. A directed diamond
has two length-two paths with the same endpoints; subtracting them cancels
the possibly missing shortcut edge. Its boundary equates the two path sums.
Vertex gradients satisfy every such relation, as does the cochain that is
one on each edge. We show that these are the only solutions with high
probability. A single transitive triangle then excludes the constant cochain.
The main difficulty is uniformity: a cochain under examination can depend
on the entire graph. Operator estimates, rather than estimates for a fixed
cochain, provide the required simultaneous conclusion.

\section{The diamond quadratic form}

In this section $G$ is a fixed loop-free digraph, with edge set
$E=\{(a,b):A_{ab}=1\}$. We work first over $\R$.
For an edge cochain $f\in\R^E$ and a vertex function $h\in\R^n$, define
$(\delta h)_{ab}=h_b-h_a$. The space left undetected by the diamond
relations will be
\begin{equation}\label{eq:W}
 W=\im\delta+\R\one_E.
\end{equation}
All vector norms and inner products are Euclidean, and all matrix norms
are their induced operator norms. The letter $J$ denotes an all-one matrix
of the dimensions specified by its indices.

A diamond consists of four distinct vertices $a,b,c,e$ and edges
$a\to b$, $b\to e$, $a\to c$, $c\to e$. It is counted once for ordered
endpoints $(a,e)$ and an unordered pair of middle vertices $\{b,c\}$.
Its boundary row and energy are
\[
 q_{a;b,c;e}=[a,b]+[b,e]-[a,c]-[c,e],\qquad
 \mathcal E_G(f)=\sum_{\text{diamonds}}\ip{q_{a;b,c;e}}{f}^{\,2}.
\]
The sign chosen for the unordered middle pair is immaterial.
The row is the literal boundary of $[a,b,e]-[a,c,e]$:
the two occurrences of $[a,e]$ cancel. Thus this two-chain belongs to
the non-regular $\Omega_2$, whether or not $a\to e$ is an edge.
Every row annihilates $W$ by telescoping and by $1+1-1-1=0$.

To expand the energy, put
\[
 O_v=\{b:A_{vb}=1\},\quad I_v=\{a:A_{av}=1\},\quad
 x_v=(f_{vb})_{b\in O_v},\quad y_v=(f_{av})_{a\in I_v}.
\]
Their sums will be denoted by
$r_v=\sum_{b\in O_v}f_{vb}$, $c_v=\sum_{a\in I_v}f_{av}$,
and $s=\sum_{(a,b)\in E}f_{ab}$. Empty vectors and sums are interpreted
as zero. Let $D_{ab}$ be the number of diamonds containing an edge $(a,b)$.
The two fork matrices have zero diagonal and, off the diagonal, entries
\begin{align}
 (P_a)_{b,c}&=\sum_{e\notin\{a,b,c\}}A_{be}A_{ce}
       &&(b,c\in O_a,\ b\ne c),\label{eq:P}\\
 (Q_e)_{b,c}&=\sum_{a\notin\{e,b,c\}}A_{ab}A_{ac}
       &&(b,c\in I_e,\ b\ne c).\label{eq:Q}
\end{align}
For an adjacent pair of edges through $v$, use the rectangular matrix
\begin{equation}\label{eq:C}
 (C_v)_{a,b}=\one_{a\ne b}
       \sum_{c\notin\{a,b,v\}}A_{ac}A_{cb}
       \quad(a\in I_v,\ b\in O_v).
\end{equation}
Finally the opposite-edge matrix, indexed by actual edges, is
\begin{equation}\label{eq:H}
 H_{(a,b),(c,e)}=
 \one_{\{a,b,c,e\}\text{ has size }4}\,A_{ac}A_{be}
       \quad((a,b),(c,e)\in E).
\end{equation}
This matrix need not be symmetric; its operator norm will suffice.

\begin{lemma}\label{lem:identity}
For every real edge cochain,
\begin{equation}\label{eq:identity}
 \mathcal E_G(f)=\sum_E D_{ab}f_{ab}^2
  -\sum_v x_v^TP_vx_v-\sum_v y_v^TQ_vy_v
  +2\sum_v y_v^TC_vx_v-2f^THf.
\end{equation}
\end{lemma}
\begin{proof}
For one diamond write its four values as $x=f_{ab}$, $y=f_{be}$,
$u=f_{ac}$, $z=f_{ce}$. Expansion of $(x+y-u-z)^2$ gives diagonal
squares, positive adjacent products $2xy+2uz$, negative source and sink
products $-2xu-2yz$, and negative opposite products $-2xz-2yu$.
The diagonal multiplicity is $D_{ab}$. Each fork matrix counts both orders
of its two middle vertices, producing its coefficient $-2$.
The matrix $C_v$ counts each adjacent product once. The two orders of
$(b,c)$ in $H$ produce the two opposite products. These are precisely
the five terms in~\eqref{eq:identity}.
\end{proof}

We will prove that~\eqref{eq:identity} is close, uniformly in $f$, to a
quadratic form depending only on $\norm f$, $r-c$ and $s$.
The needed estimates account for the dependence between a selected edge
and its neighboring diamond relations.

\section{Concentration ingredients}

Fix a sequence satisfying~\eqref{eq:condition}, and abbreviate
\begin{equation}\label{eq:scales}
 \ell=\log(2n),\qquad d=n^2p^3,\qquad \mu=(n-3)p^2,
 \qquad
 \varepsilon_n=\sqrt d\,\ell^2
     +d\sqrt{\frac{\ell}{np}}+np^2\ell^2+\ell^3.
\end{equation}
In particular $d\to\infty$ and $np=(nd)^{1/3}$ grows faster than every
fixed power of $\ell$. Constants denoted $C_L$ depend only on a fixed
positive number $L$ and may increase finitely many times. A failure bound
$O_L(n^{-L})$ has this same uniformity in $n$, $p$ and the fixed vertex
sets under consideration.

The only external concentration input is matrix Bernstein in the following
two forms~\cite[Theorems 1.4 and 1.6 of the author preprint]{Tropp}.
For independent mean-zero self-adjoint $q\times q$ matrices $X_i$ with
$\norm{X_i}\leq R$, put $\sigma^2=\norm{\sum_i\E X_i^2}$. Then
\begin{equation}\label{eq:bernstein}
 \Prob\left\{\norm{\sum_iX_i}\geq t\right\}
 \leq2q\exp\left\{-\frac{t^2}{2(\sigma^2+Rt/3)}\right\}\qquad(t\geq0).
\end{equation}
The two-sided form follows by applying the source's upper-eigenvalue
bound to both $X_i$ and $-X_i$.
For independent mean-zero $q_1\times q_2$ matrices $Z_i$ with
$\norm{Z_i}\leq R$, the norm bound has prefactor $q_1+q_2$ and the
same exponent, with
\[
 \sigma^2=\max\left\{\norm{\sum_i\E Z_iZ_i^T},
                          \norm{\sum_i\E Z_i^TZ_i}\right\}.
\]
Taking dimension one gives, for a sum of independent Bernoulli variables
with mean $M$, deviation at most $C_L(\sqrt{M\ell}+\ell)$ with failure
$O_L(n^{-L})$. Throughout the proof dimensions are at most $n^2$, so
the logarithmic factor absorbs the dimension prefactors and polynomial
union bounds by choosing $L$ larger at the outset.

\subsection{Independent row and column restrictions}

The next lemma gains a factor $p$ when both coordinate sets are sampled.
It will be used only where the sampled edge families are independent of
the matrix being restricted.

\begin{lemma}\label{lem:restriction}
Let $M$ be a deterministic real matrix with both dimensions at most $n^2$.
Let $R$ and $S$ be independent diagonal matrices whose diagonal entries
are independent Bernoulli$(p)$ variables. If
\[
 a=\norm M,\quad u=\max_i\norm{M_{i,*}}_2,\quad
 v=\max_j\norm{M_{*,j}}_2,\quad b=\max_{i,j}|M_{ij}|,
\]
then, with failure $O_L(n^{-L})$,
\begin{equation}\label{eq:restriction}
 \norm{RMS}\leq pa+C_L\sqrt{p\ell}(u+v)+C_L\ell b.
\end{equation}
The same conclusion holds conditionally for a random $M$ independent of
$R$ and $S$.
\end{lemma}
\begin{proof}
First consider row selection of a fixed matrix $B$ with rows $b_i$.
Writing $u_B=\max_i\norm{b_i}_2$, the centered Gram matrix is
\[
 B^TRB-pB^TB=\sum_i(\xi_i-p)b_i^Tb_i.
\]
Its independent summands have norm at most $u_B^2$. Since
$(b_i^Tb_i)^2=\norm{b_i}_2^2 b_i^Tb_i$, their variance matrix has norm
at most $p u_B^2\norm B^2$. The upper-eigenvalue bound in
\eqref{eq:bernstein} gives
\[
 \norm{RB}^2\leq p\norm B^2
       +C_L u_B\sqrt{p\ell}\norm B+C_L\ell u_B^2.
\]
Taking square roots and enlarging $C_L$ yields
\begin{equation}\label{eq:one-sided}
 \norm{RB}\leq\sqrt p\norm B+C_L\sqrt\ell\,u_B.
\end{equation}
Apply~\eqref{eq:one-sided} to $M^T$ to obtain
$\norm{MS}\leq\sqrt p\,a+C_L\sqrt\ell\,v$.
For a fixed row, $\norm{(MS)_{i,*}}_2^2=\sum_j\zeta_j|M_{ij}|^2$.
The centered summands are bounded by $b^2$ and have total variance at most
$pb^2u^2$. Scalar Bernstein and a union bound over the rows give
\[
 \max_i\norm{(MS)_{i,*}}_2\leq\sqrt p\,u+C_L\sqrt\ell\,b.
\]
Conditional on $S$, apply~\eqref{eq:one-sided} to $MS$ using independent
$R$. Substituting the preceding two inequalities proves
\eqref{eq:restriction}. Conditioning on an independent random $M$
proves the last assertion.
\end{proof}

\subsection{Products on disjoint vertex blocks}

For disjoint vertex sets $V_i$, write $B_{ij}=(A_{ab})_{a\in V_i,b\in V_j}$.
The entries of each block are independent Bernoulli$(p)$; blocks with
different ordered source--target sets are independent.

\begin{lemma}\label{lem:blocks}
Fix three disjoint vertex sets of sizes at most $n$, and put
$M=B_{12}B_{23}-|V_2|p^2J$. For independent Bernoulli$(p)$ coordinate
selectors $R,S$, also independent of the blocks,
\begin{equation}\label{eq:product-bound}
 \norm{RMS}\leq C_L\varepsilon_n
\end{equation}
with failure $O_L(n^{-L})$.
For four disjoint vertex sets the same bound holds with
$M=B_{13}\otimes B_{24}-p^2J$, where the row coordinates are pairs
in $V_1\times V_2$ and the column coordinates are pairs in
$V_3\times V_4$.
\end{lemma}
\begin{proof}
For any one of these blocks let $X=B-pJ$. Entrywise rectangular
Bernstein has summand norm at most $1$ and variance at most $np$, giving
\begin{equation}\label{eq:block-spectral}
 \norm X\leq C_L\sqrt{np\ell}
\end{equation}
because $np\gg\ell$. Scalar Bernstein and union bounds give degrees at
most $C_Lnp$ and absolute centered row and column sums at most
$C_L\sqrt{np\ell}$. Each row and column of $B$ and $X$ has norm at most
$C_L\sqrt{np}$: for $X$, use
$\sum_j(B_{ij}-p)^2\leq2\sum_jB_{ij}+2np^2$.
These bounds are uniform in every fixed choice of the vertex sets,
including empty sets.

For three sets, the centered product is
\[
 M=pJX_{23}+pX_{12}J+X_{12}X_{23}.
\]
Using~\eqref{eq:block-spectral} and $\norm J\leq n$ gives
\begin{equation}\label{eq:full-product}
 \norm M\leq C_L\bigl(np\ell+n^{3/2}p^{3/2}\sqrt\ell\bigr).
\end{equation}
For each row the three terms are bounded respectively by
$p\sqrt n\norm{X_{23}}$, $p\sqrt n$ times the magnitude of a row
sum of $X_{12}$, and a row norm of $X_{12}$ times
$\norm{X_{23}}$. Columns satisfy the reversed estimates. Thus
\begin{equation}\label{eq:product-rows}
 \max\{\text{row norms, column norms of }M\}
       \leq C_Lnp\sqrt\ell.
\end{equation}
Each entry of $B_{12}B_{23}$ is Binomial$(|V_2|,p^2)$: different middle
vertices use distinct edges of the two disjoint blocks. Consequently
\begin{equation}\label{eq:product-entries}
 \max|M_{ij}|\leq C_L\bigl(\sqrt{np^2\ell}+\ell\bigr)
                   \leq C_L(np^2+\ell).
\end{equation}
Lemma~\ref{lem:restriction} applied to
\eqref{eq:full-product}--\eqref{eq:product-entries} now yields
\[
 \norm{RMS}\leq C_L\left(
 np^2\ell+n^{3/2}p^{5/2}\sqrt\ell+\sqrt d\,\ell+\ell^2\right)
 \leq C_L\varepsilon_n.
\]
Here $n^{3/2}p^{5/2}\sqrt\ell=d\sqrt{\ell/(np)}$.

For the tensor product expand instead
\[
 M=X_{13}\otimes X_{24}+pJ\otimes X_{24}+pX_{13}\otimes J.
\]
The inequality $\norm{U\otimes V}\leq\norm U\norm V$ gives
\eqref{eq:full-product} again. The uncentered tensor has row and column
norms bounded by products of two block row or column norms, at most
$C_Lnp$. A row or column of the subtracted matrix $p^2J$ has norm
at most $np^2\leq np$. Entry magnitudes of $M$ are at most $2$.
Lemma~\ref{lem:restriction} therefore gives
\[
 \norm{RMS}\leq C_L\left(
 np^2\ell+n^{3/2}p^{5/2}\sqrt\ell+\sqrt{d\ell}+\ell\right)
 \leq C_L\varepsilon_n.
\]
The tensor norm inequality follows, for example, by reshaping a vector
as a matrix and using $\norm{UZV}_F\leq\norm U\norm Z_F\norm V$.
\end{proof}

\subsection{Averaging the cuts}

We use random vertex labels to place the selected edges and connector edges
in independent blocks. The following finite averaging observation recovers
the original matrix without demanding that every partition be good.

\begin{lemma}\label{lem:averaging}
Let $k,t$ be fixed positive integers. Suppose a graph-dependent matrix $T$
and matrices $T_\chi$ associated with independent uniform $k$-labels on
vertices satisfy $T=k^t\E_\chi T_\chi$. The matrices are embedded in a
common coordinate space for each fixed graph. Suppose every
$\norm{T_\chi}\leq2n^2$, and for each fixed labeling
\[
 \Prob_G\{\norm{T_\chi}>C_L\varepsilon_n\}\leq O_L(n^{-L}).
\]
Then, except with probability $O_L(n^{-L/2})$,
\[
 \norm T\leq k^t\bigl(C_L\varepsilon_n+2n^{2-L/2}\bigr).
\]
\end{lemma}
\begin{proof}
Let $h(G)$ be the fraction of labelings with
$\norm{T_\chi}>C_L\varepsilon_n$. Finite-sum interchange gives
$\E_G h(G)\leq O_L(n^{-L})$.
By Markov's inequality, $h(G)\leq n^{-L/2}$ outside an event of
probability $O_L(n^{-L/2})$. On that complement the triangle inequality
for the finite average gives the stated bound.
\end{proof}

\section{Uniform diamond estimates}

We now return to the matrices~\eqref{eq:P}--\eqref{eq:H} indexed by the
random neighborhoods or edge set of $G(n,p)$.

\begin{proposition}\label{prop:matrices}
Under~\eqref{eq:condition}, with probability tending to one, the following
bounds hold simultaneously:
\begin{align}
 \max_{(a,b)\in E}|D_{ab}-2d|&\leq C\varepsilon_n,\label{eq:D-bound}\\
 \max_a\norm{P_a-\mu(J-I)}+
 \max_e\norm{Q_e-\mu(J-I)}&\leq C\varepsilon_n,\label{eq:fork-bound}\\
 \max_v\norm{C_v-\mu J}&\leq C\varepsilon_n,\label{eq:C-bound}\\
 \norm{H-p^2J_{E,E}}&\leq C\varepsilon_n.\label{eq:H-bound}
\end{align}
Here $C$ is an absolute constant.
\end{proposition}

We prove the mixed and opposite bounds first, then the fork and diagonal
bounds. Take, for example, $L=100$ in the concentration ingredients;
enlarging initial tail exponents accommodates all finite union bounds.

\subsection{Mixed terms and opposite edges}

Fix a vertex $v$ and assign the other vertices independent uniform labels
from $\{1,2,3\}$. Let $V_i$ be their label classes. Retain in $C_v$
only triples $(a,c,b)$ of labels $(1,2,3)$, and subtract
$|V_2|p^2$ in the retained endpoint rectangle. This centered cut is the
restriction
\[
 R_{I_v}\bigl(B_{12}B_{23}-|V_2|p^2J\bigr)R_{O_v}.
\]
The row selectors are $A_{av}$ for $a\in V_1$, the column selectors
are $A_{vb}$ for $b\in V_3$. These two families are mutually independent
and independent of $B_{12},B_{23}$, since the latter blocks avoid $v$.
Lemma~\ref{lem:blocks} applies for each fixed labeling.

For a fixed graph, pad this cut with zero rows and columns into the
coordinate space $I_v\times O_v$ and average. Each distinct triple
$(a,c,b)$ has the prescribed labels with probability $1/27$.
Consequently $27$ times the average centered cut is exactly
\[
 C_v-\mu(J-F_v),
\]
where $(F_v)_{a,b}=\one_{a=b}$ on $I_v\times O_v$.
Indeed, for $a\ne b$ there are $n-3$ eligible middle vertices; for
$a=b$ the label requirements are impossible. The deterministic cut norm
is at most $2n^2$, because its dimensions are at most $n$ and its entry
magnitudes at most $2n$. Lemma~\ref{lem:averaging} and a union bound over
$v$ yield $\max_v\norm{C_v-\mu(J-F_v)}\leq C\varepsilon_n$.
Since $\norm{F_v}\leq1$ and $\mu\leq np^2\leq\varepsilon_n$,
this proves~\eqref{eq:C-bound}.

For $H$ use four labels, assigning $(a,b,c,e)$ the labels $(1,2,3,4)$.
The connector matrix is $B_{13}\otimes B_{24}$, indexed by pairs
$(a,b)$ and $(c,e)$. Restriction to actual edge pairs uses independent
selectors $A_{ab}$ in block $12$ and $A_{ce}$ in block $34$.
They are mutually independent and independent of connector blocks
$13$ and $24$. Thus Lemma~\ref{lem:blocks} applies to the centered cut
$B_{13}\otimes B_{24}-p^2J$ with these pair-coordinate selectors.
After zero-padding in the edge coordinate space, each four-distinct
tuple occurs with probability $1/256$ and the cut norm is at most $2n^2$.
Lemma~\ref{lem:averaging} gives
\[
 \norm{H-p^2J^{(4)}_{E,E}}\leq C\varepsilon_n,
\]
where $J^{(4)}$ is one for edge pairs with four distinct vertices and
zero otherwise.

With probability $1-O_L(n^{-L})$, all in- and out-degrees are at most
$C_Lnp$, by scalar Bernstein. On this event each actual edge shares
a vertex with at most $C_Lnp$ other actual edges. Hence the symmetric
nonnegative matrix $J_{E,E}-J^{(4)}_{E,E}$ has maximum row sum at most
$C_Lnp$, and its norm has the same upper bound. Multiplication by $p^2$
costs at most $C_Lnp^3\leq C_L\varepsilon_n$, proving
\eqref{eq:H-bound}. For completeness, a matrix with entry magnitudes
having row sums at most $R$ and column sums at most $S$ has norm at most
$\sqrt{RS}$, by weighted Cauchy--Schwarz. This justifies the row-sum
bound without a symmetry assumption on $H$ itself.

\subsection{Forks}

Fix $a$ and condition only on its outgoing neighborhood $O_a=S$.
On the degree event, $m=|S|\leq C_Lnp$. For $e\ne a$ let
$u_e=(A_{be})_{b\in S}$ and
$Z_e=u_eu_e^T-\diag(u_e)$. These columns are independent:
their edges have distinct ordered target vertices and none has source $a$.
When $e\in S$ its self-loop coordinate is deterministically zero; the
remaining coordinates are independent Bernoulli$(p)$.
Reciprocal edges in different columns remain independent. We have
\[
 P_a=\sum_{e\ne a}Z_e,\qquad \E(P_a\mid S)=\mu(J-I).
\]
For each off-diagonal pair $(b,c)$ the $n-3$ permissible witnesses exclude
$a,b,c$, which explains the exact mean.

For one such column, the diagonal entries of $\E Z_e^2$ are at most
$mp^2$, and the off-diagonal entries are at most $mp^3$.
This follows from $(Z_e)_{ij}=(u_e)_i(u_e)_j$ for $i\ne j$:
a diagonal square uses two distinct Bernoulli coordinates, while an
off-diagonal square uses three. The entries are nonnegative, so their
maximum row sum is at most $mp^2+m^2p^3$.

Truncate each $Z_e$ separately to zero if the column support satisfies
$|\supp u_e|>C_L(mp+\ell)$, and call the result $\widetilde Z_e$.
Scalar Bernstein bounds the probability of each truncation by
$O_L(n^{-L})$. Truncation preserves column independence, and the total
mean lost has norm at most $O_L(nm\,n^{-L})$, since $\norm{Z_e}\leq m$.
Every entry of $\E\widetilde Z_e^2$ is bounded by the corresponding
nonnegative entry of $\E Z_e^2$. Thus its norm is at most
$mp^2+m^2p^3$, by the same row-sum bound. In addition,
\[
 \E(\widetilde Z_e-\E\widetilde Z_e)^2
 =\E\widetilde Z_e^2-(\E\widetilde Z_e)^2
 \preccurlyeq\E\widetilde Z_e^2.
\]
The centered variance sum therefore has norm at most
\[
 n(mp^2+m^2p^3)\leq C_Ld(1+np^2).
\]
Each centered summand has norm at most $C_L(np^2+\ell)$, because the
retained support count bounds $\norm{\widetilde Z_e}$ and
$mp\leq C_Lnp^2$. Applying~\eqref{eq:bernstein} gives
\begin{equation}\label{eq:fork-deviation}
 \norm{P_a-\mu(J-I)}
 \leq C_L\left(\sqrt{d(1+np^2)\ell}+(np^2+\ell)\ell\right)+o(1)
 \leq C_L\varepsilon_n.
\end{equation}
Here $\sqrt{d\,np^2\ell}=d\sqrt{\ell/(np)}$.
The truncated and untruncated sums differ only on the union of the
column truncation events. The conditional estimates are uniform in
every $S$ of the permitted size, so they may be averaged over $S$ and
union bounded over $a$. Empty or singleton neighborhoods cause no
difficulty. Reversing all arrows gives the corresponding bound for
$Q_e$. This proves~\eqref{eq:fork-bound}.

\subsection{Diagonal multiplicities}

Fix $(a,b)$ and condition on $A_{ab}=1$. The diamonds in which this edge
is on the first step are counted by
\[
 \sum_{\substack{c,e\notin\{a,b\}\\c\ne e}}
             A_{ac}A_{be}A_{ce}.
\]
Expose just the outgoing rows of $a,b$, and put
$S=O_a\setminus\{b\}$ and $T=O_b\setminus\{a\}$.
The remaining edges $c\to e$ with $(c,e)\in S\times T$, $c\ne e$,
have source different from $a,b$. They are unexposed independent
Bernoulli$(p)$ variables. The conditional count is consequently
\[
 \operatorname{Binomial}\bigl(|S||T|-|S\cap T|,p\bigr).
\]
Under $A_{ab}=1$, scalar Bernstein for the other row entries gives
\[
 |S|,|T|=np+O_L(\sqrt{np\ell}+\ell+1)
\]
outside an event of probability $O_L(n^{-L})$.
Since $|S\cap T|\leq\min(|S|,|T|)$ and $np\gg\ell$, the conditional
mean differs from $d$ by at most $C_Ld\sqrt{\ell/(np)}$.
Another conditional scalar Bernstein bound contributes at most
$C_L(\sqrt{d\ell}+\ell)$. The first-step count is therefore
$d+O_L(\varepsilon_n)$.

For diamonds in which $(a,b)$ is on the second step, expose instead the
incoming edge families of $a,b$. The remaining connectors have target
different from $a,b$, so the same conditional calculation gives
$d+O_L(\varepsilon_n)$. These two types are disjoint and their sum is
$D_{ab}$; they need not be independent. The failure bound is uniform
conditional on every possible present edge, and a union bound over
at most $n^2$ edges proves~\eqref{eq:D-bound}.
No conditioning on a global degree event is added to the conditional
binomial law. This finishes the proof of Proposition~\ref{prop:matrices}.

\subsection{The uniform quadratic-form estimate}

The estimates just proved hold on a single graph event and therefore
apply simultaneously to graph-dependent cochains.

\begin{proposition}\label{prop:energy}
With probability tending to one under~\eqref{eq:condition}, every
$f\in\R^E$ satisfies
\begin{equation}\label{eq:energy-bound}
 \left|\mathcal E_G(f)-
 \bigl((2d+2\mu)\norm f^2-\mu\norm{r-c}^2-2p^2s^2\bigr)\right|
 \leq C\varepsilon_n\norm f^2.
\end{equation}
In particular, on that event and for all sufficiently large $n$,
\begin{equation}\label{eq:coercivity}
 f\perp W\quad\Longrightarrow\quad
          \mathcal E_G(f)\geq d\norm f^2.
\end{equation}
\end{proposition}
\begin{proof}
In~\eqref{eq:identity}, replacing the five matrix families by their main
terms in Proposition~\ref{prop:matrices} gives respectively
\[
 2d\norm f^2,\quad
 -\mu(\norm r^2-\norm f^2),\quad
 -\mu(\norm c^2-\norm f^2),\quad
 2\mu\ip c r,\quad -2p^2s^2.
\]
Their sum is the main form in~\eqref{eq:energy-bound}.
For each fork family, the error is at most $C\varepsilon_n\norm f^2$
because every edge appears in exactly one outgoing vector and exactly
one incoming vector. For the mixed term use
\[
 \sum_v\norm{x_v}\norm{y_v}
 \leq\left(\sum_v\norm{x_v}^2\right)^{1/2}
      \left(\sum_v\norm{y_v}^2\right)^{1/2}=\norm f^2.
\]
The diagonal and opposite-edge errors have the same bound by
\eqref{eq:D-bound} and~\eqref{eq:H-bound}. This proves
\eqref{eq:energy-bound} uniformly in $f$.

For $f\perp W$, pairing with gradients gives
$\ip f{\delta h}=\sum_vh_v(c_v-r_v)=0$ for every $h$, so $r=c$.
Pairing with $\one_E$ gives $s=0$. Finally
\[
 \frac{\varepsilon_n}{d}
 =\frac{\ell^2}{\sqrt d}+\sqrt{\frac{\ell}{np}}
        +\frac{\ell^2}{np}+\frac{\ell^3}{d}\longrightarrow0
\]
by~\eqref{eq:condition} and $np=(nd)^{1/3}$.
Thus $(2d+2\mu)-C\varepsilon_n\geq d$ for all sufficiently large $n$,
which proves~\eqref{eq:coercivity}.
\end{proof}

\section[Rational homology vanishes]{\hspace{0.6em}Rational homology vanishes}

We complete the proof of Theorem~\ref{thm:main}.
Let $K$ be the space of real edge cochains annihilating every diamond
row. We already have $W\subseteq K$. On the event of
Proposition~\ref{prop:energy}, any $f\in K\setminus W$ would have a
nonzero orthogonal component $f_0\in W^\perp$. Since subtracting a
member of $W$ preserves all diamond evaluations, $f_0\in K$ and
$\mathcal E_G(f_0)=0$. This contradicts~\eqref{eq:coercivity}.
Consequently
\begin{equation}\label{eq:annihilator}
 K=\im\delta+\R\one_E
\end{equation}
with probability tending to one.

The graph also contains a transitive triangle with probability tending
to one. To see this, fix $v=1$. Scalar Bernstein gives
$|O_v|=(1+o(1))np$ with high probability. Conditional on $O_v$, all
ordered edges between distinct vertices of $O_v$ are independent
Bernoulli$(p)$: none has source $v$. Their probability of all being absent
is at most
\[
 (1-p)^{|O_v|(|O_v|-1)}
 \leq \exp\{-p|O_v|(|O_v|-1)\}
 =\exp\{-(1+o(1))d\}\longrightarrow0.
\]
An edge $a\to b$ inside $O_v$ supplies $v\to a$, $a\to b$, $v\to b$.
The allowed two-path $[v,a,b]$ has allowed boundary
\[
 [a,b]-[v,b]+[v,a].
\]
Every gradient evaluates to zero on this row, whereas $\one_E$ evaluates
to $1$. Combining this triangle with~\eqref{eq:annihilator}, the common
annihilator of the diamond and transitive-triangle rows is exactly
$\im\delta$.

Let $\mathcal B\subseteq\R^E$ be the span of these rows and let
$\partial_1:\R^E\to\R^n$ be the vertex-incidence boundary.
All the rows are cycles, and
\[
 \mathcal B^\perp=\im\delta=(\ker\partial_1)^\perp.
\]
Finite-dimensional duality yields $\mathcal B=\ker\partial_1$.
The incidence matrix and the matrix of the displayed boundary rows have
integer entries, so each has the same rank over $\Q$ and $\R$:
rank is the largest order of a nonzero minor. Therefore the rational span
of these valid non-regular boundaries is the full rational edge-cycle
space. It is contained in $\partial\Omega_2$, which is itself contained
in $\ker\partial_1$. Their equality proves
$\beta_1^{\rm nr}(G(n,p_n))=0$ on an event of probability tending to one,
as required.

\paragraph{AI assistance.}
OpenAI Codex was used substantially for proof exploration, manuscript
drafting, and automated mathematical checking. Automated checks do not
constitute independent human peer review. The paper supplies the full
argument and its external concentration input for further verification.

\begin{thebibliography}{9}
\bibitem{Chaplin}
T. Chaplin,
\emph{First Betti number of the path homology of random directed graphs},
Journal of Applied and Computational Topology \textbf{8} (2024), 1503--1549.
\href{https://doi.org/10.1007/s41468-022-00108-3}{doi:10.1007/s41468-022-00108-3}.

\bibitem{Tropp}
J. A. Tropp,
\emph{User-Friendly Tail Bounds for Sums of Random Matrices},
Foundations of Computational Mathematics \textbf{12} (2012), 389--434.
\href{https://doi.org/10.1007/s10208-011-9099-z}{doi:10.1007/s10208-011-9099-z}.
Author preprint, revised June 15, 2011,
\href{https://tropp.caltech.edu/papers/Tro11-User-Friendly-preprint.pdf}{Theorems 1.4 and 1.6}.
\end{thebibliography}
\end{document}
